\section{Discussion: Toward VLM-Native Media}

The experiments above are intentionally conservative: frozen models, explicit
evidence classes, no speed claim where the backend did not skip timed work.
That conservatism should not hide the bigger implication. Existing media
pipelines already contain cheap temporal structure, but most video VLM stacks
consume the final RGB frames as if every frame were equally new. Anti-
recomputation asks a systems question before it asks for a new model: what did
the media pipeline already know, and why did the VLM pay to rediscover it?

\subsection{Unchanged Is Only The First Case}

Same-position reuse is the simplest regime: if a patch did not change, reuse
its visual evidence. But video codecs also exploit a richer fact: many things
change predictably. Objects translate, cameras pan, forward-driving cameras
produce structured expansion, and occlusions create localized new evidence.
The current paper earns evidence mostly for static-to-medium-motion benchmark
regimes and for same-video follow-up reuse. It does \emph{not} yet prove
motion-compensated feature reuse, driving-specific routing, or sensor-fusion
caching. Those directions follow naturally from the failures and ceilings we
measured, so they belong in the paper's agenda rather than in the headline
claim.

\begin{table}[H]
\centering
\caption{Application regimes as evidence targets, not generic motivation. The
right policy depends on the redundancy structure and on what failure would
cost.}
\label{tab:application-regimes}
\footnotesize
\renewcommand{\arraystretch}{1.22}
\begin{tabularx}{\linewidth}{@{}
  >{\raggedright\arraybackslash}p{0.18\linewidth}
  >{\raggedright\arraybackslash}p{0.22\linewidth}
  >{\raggedright\arraybackslash}p{0.26\linewidth}
  >{\raggedright\arraybackslash}X
@{}}
\toprule
Regime & Redundancy structure & Candidate routing policy & Evidence status / risk \\
\midrule
Static cameras, factories, surveillance, stove/countertop & stable background,
localized entrants, rare decisive events & persistent visual memory plus
bounded-staleness refresh around moving regions & closest to current evidence;
needs stronger application baselines and event-window recall \\
\addlinespace[0.25em]
Screen and UI video & exact-copy regions, glyph changes, scrolling, cursor
motion & screen-content path with exact-copy, OCR/glyph, and high-frequency
guards & planned specialization; current VideoMME audit under-samples scroll/pan \\
\addlinespace[0.25em]
Forward driving & structured ego-motion, border entry, independently moving
objects & pose/flow-conditioned refresh with protected object and boundary
regions & future lane; should not be inferred from same-position reuse \\
\addlinespace[0.25em]
Egocentric, FPV, drones & jerky motion, parallax, frequent occlusion & global
stabilization or multi-reference cache before reuse & current benchmark mix is
the wrong corpus for this claim \\
\addlinespace[0.25em]
Robotics / VLA & static workspace plus gripper/object/contact zones and hard
deadlines & protected ROIs, p95/p99 latency budgets, guarded dense fallback & future
systems lane; answer-stability QA is not a control-safety metric \\
\addlinespace[0.25em]
Games and HUD-heavy streams & stable HUD, structured camera regimes, repeated
interaction states & HUD anchors plus event-triggered world refresh & future
application lane; likely high reuse but unmeasured here \\
\bottomrule
\end{tabularx}
\end{table}

\paragraph{Robotics is a different contract.}
Robotics is not just another row in an application table. It changes the
optimization target from average answer stability to deadline-safe perception.
A plausible safety contract is protected reuse, not maximal reuse: test
policies that always refresh the gripper, manipulated object, contact boundary,
goal receptacle, newly entering regions, and any sensor-disagreement zone;
reuse aggressively elsewhere; and report p95/p99 latency, deadline misses,
fallback rate, and task success instead of only QA accuracy. We do not claim
that evidence here. We name it because it is the natural systems lane once
anti-recomputation leaves offline QA.

\subsection{From Video Streams To State Streams}

The scale-out streaming rows are the concrete predecessor of this interface:
they are not yet a standardized C-STREAM result, but they show why a
perception-rate state stream is the right deployment abstraction to test.

The same logic extends beyond RGB video. A robot, vehicle, or factory cell may
carry multiple cameras, depth or ToF, IMU/odometry, event sensors, object
tracks, calibration, timestamps, and confidence estimates. These side channels
are not semantic-saliency oracles. They are cheaper refresh and uncertainty
oracles: RGB says what a region looks like; depth says whether geometry moved;
IMU says whether apparent motion is explained by the camera; event sensors
flag onset; multi-camera geometry distinguishes occlusion from disappearance.
The VLM-native interface we want is therefore not ``more frames,'' but a
synchronized state-update stream that tells the runtime where fresh visual
evidence is worth buying.

\subsection{Codec Signals Are Requirements Probes}

Classical codecs are hand-engineered prediction systems: predict a reference,
send motion and residual surprise, transform and quantize under a rate-
distortion objective, and make the result cheap enough for hardware. That
objective was designed mainly for human viewing. Standards work is already
moving toward machine consumption, including MPEG-AI Part 2 Video Coding for
Machines \cite{mpegvcm}, JPEG AI \cite{jpegai}, and event-data coding efforts
such as JPEG XE \cite{jpegxe}. Older object/layer media ideas, including
MPEG-4 Visual object coding \cite{mpeg4visual}, also show that rectangular RGB
frames were not the only conceivable abstraction. The analogy to object-based
audio such as Dolby Atmos \cite{dolbyatmos} is useful but only as an analogy:
the goal is not theatrical rendering, but a synchronized stream of
model-facing state updates.

Our role in that trajectory is narrower. We are not designing the future codec
here. We use VLM behavior to motivate candidate requirements for that future
interface:

\begin{table}[H]
\centering
\caption{How current limitations motivate future media-interface requirements.
These are hypotheses and candidate design requirements, not current
measured claims.}
\label{tab:media-requirements}
\small
\renewcommand{\arraystretch}{1.14}
\begin{tabularx}{\linewidth}{@{}
  >{\raggedright\arraybackslash}p{0.40\linewidth}
  >{\raggedright\arraybackslash}X
@{}}
\toprule
Observed limitation or boundary & Future VLM-native media should expose \\
\midrule
Mean block difference misses small but decisive events & residual
concentration, changed-pixel fraction, onset flags, and max-age state \\
Novelty-ranked dense frames can cluster budget in the wrong place & temporal
coverage constraints, event-window hints, and placement-aware scheduling \\
Same-position reuse weakens under camera motion & camera pose, global motion,
depth-aware flow, and boundary-entry flags \\
Predictably moving objects are not static but are not arbitrary either &
object IDs, masks, tracks, motion uncertainty, and multi-reference state \\
Human-compression artifacts can hide model-critical text or color & task-aware
text, chroma, palette, and high-frequency metadata \\
Sparse regions may be awkward for today's tensor kernels & tensor-friendly
active tiles, projector-consistent groups, and regular batched sparse layouts \\
RGB alone confuses lighting, geometry, and contact & synchronized depth/ToF,
event, IMU/odometry, timestamp, and confidence sidecars \\
\bottomrule
\end{tabularx}
\end{table}

\paragraph{Layered world-state media.}
One useful mental model is a machine-facing analogue of object/layer media:
static background memory, camera pose, object tracks and masks, motion
uncertainty, residual surprise, text/glyph events, event-camera onset,
depth/ToF changes, timestamps, and confidence. The model would render or
refresh dense visual evidence only where the state update requires it. This is
not a current codec implementation; it is the interface pressure revealed by a
paper about paying for the wall again.

\paragraph{Beyond zero-change inference.}
The current result is useful precisely because it is frozen-model and
training-free. It identifies where an existing model already tolerates reuse
and where it does not. The model-changing version should train around those
traces rather than ignore them: learned recache gates, codec-conditioned token
types, P-frame or delta encoders, and cache-state-aware attention can all use
the failure cases and routing logs produced by this paper as supervision.
There is also a training-compute hypothesis: if future VLMs learn from
state-update streams rather than dense RGB frame stacks, they may spend the
same training FLOPs on more temporal coverage, more queries per clip, or harder
event windows. We do not measure that here. The caveat is important:
post-vision-tower feature caching does not directly become an end-to-end training
speedup when the vision encoder is trainable, because cached features
approximate or break the gradient path. The training version likely needs
model-native delta encoders, differentiable refresh gates, or teacher/student
traces rather than dropping inference caches into pretraining unchanged.

\paragraph{Beyond video.}
The general pattern is not unique to pixels. Any expensive encoder over a
slowly changing state stream can, in principle, benefit from a cheap freshness
oracle: financial order books, industrial telemetry, scientific sensor arrays,
long medical monitoring streams, or logs from embodied agents. Video is the
right starting point because temporal redundancy is explicit, codecs already
extract motion and residual signals, and the cost of visual encoding is large
enough to matter. Other domains need their own state variables and failure
contracts; the transferable lesson is anti-recomputation under a measured
quality--compute frontier, not pixel differencing itself. The transfer rule is
strict: the domain needs slowly changing state, an oracle cheaper than the
encoder, expensive recomputation, and task-specific paired-failure metrics.
Without all four, anti-recomputation is just a metaphor.

\subsection{Open Experiments}

Several high-value brainstorm items remain deliberately outside the current
evidence boundary.

The near-term research questions are reader-facing, not bookkeeping. Does
semantic substitution become a real wall-clock win when the backend actually
skips measured vision work? Does temporal placement of fresh compute beat
novelty magnitude under matched budgets? Can saved compute buy broader
temporal coverage rather than only lower latency? Can answer-margin and
paired-drift features predict when reuse will fail? Which application baselines
make deployment claims credible rather than anecdotal? How far does paired
stability extend when a user asks dialogue-like follow-up questions about the
same video, and what repair schedule prevents drift from accumulating?

The longer-range agenda is the VLM-native media ladder: object/state delta
sidecars; sensor-fusion world-state streams with depth/ToF, event, pose, and
timestamp layers; compute-denial robustness; learned recache gates;
codec-conditioned token types; learned delta encoders; and hardware-aware
active-tile decode. These ideas follow from the measurements above, but they
remain future work unless they receive their own reproduced artifacts.

The same anti-recomputation principle may also rotate into generative video
and audio. A video generator or world model could reuse latent hidden state
across denoising steps or autoregressive chunks when predicted motion and
residual uncertainty are low, then refresh around occlusion, entropy spikes, or
new-object events. Audio has a similar shape but a different trap: quiet is not
the same as irrelevant, so onset, phoneme boundaries, speaker changes, and
prosody matter more than energy alone. These are domain rotations, not evidence
claims in this paper.

The long-term version is simple to state: video for machines should become a
timeline of state updates, not only dense RGB frames. Static background,
camera motion, dynamic objects, text, events, depth, sensor confidence, and
time can all be first-class evidence. This paper is the first rung: a careful
training-free measurement of where anti-recomputation works today, where it
fails, and what future media systems should stop making the model rediscover.
